\section{System Design}
\label{sec:design}

This section describes the architecture and design decisions for SecureFL-IDS, building upon the baseline PP-FL-DP-IDS framework while incorporating enhancements for adaptive privacy and communication efficiency.

\subsection{System Architecture}

Figure~\ref{fig:architecture} illustrates the overall system architecture consisting of federated clients (cloud tenants), a central aggregator server, and data preprocessing pipelines.

\subsubsection{Federated Client}

Each client represents a cloud tenant with local network traffic data. Key components include:

\begin{itemize}
    \item \textbf{Data Loader:} Preprocesses raw network traffic into feature vectors
    \item \textbf{Local Model:} CNN or CNN-LSTM classifier for intrusion detection
    \item \textbf{Local Trainer:} Trains model on client data for specified epochs
    \item \textbf{Privacy Mechanism:} Applies differential privacy to model updates before transmission
    \item \textbf{Communication Module:} Sends privatized updates to server, receives global model
\end{itemize}

\subsubsection{Aggregation Server}

The server coordinates federated learning without accessing raw client data:

\begin{itemize}
    \item \textbf{Aggregator:} Combines client updates using FedAvg or quality-weighted averaging
    \item \textbf{Compression Module:} (Improved only) Applies top-k sparsification to reduce bandwidth
    \item \textbf{Global Model:} Maintains and broadcasts updated global parameters
    \item \textbf{Evaluation Module:} Assesses global model on held-out test set
\end{itemize}

\subsection{Data Pipeline}

\subsubsection{Preprocessing}

Raw UNSW-NB15 data undergoes:

\begin{enumerate}
    \item Missing value imputation (zero-fill for numeric features)
    \item Categorical encoding (one-hot or label encoding)
    \item Feature selection (retain 20 most informative features)
    \item Normalization (StandardScaler: zero mean, unit variance)
    \item Train/test split (80\%/20\%, stratified by label)
\end{enumerate}

\subsubsection{Non-IID Distribution}

Training data partitions across clients via Dirichlet($\alpha=0.5$) sampling, ensuring each client receives heterogeneous attack type distributions representing different tenant security profiles.

\subsection{Model Architectures}

\subsubsection{Baseline: CNN Classifier}

\begin{verbatim}
Input (batch, 20 features)
  ↓ Unsqueeze → (batch, 1, 20)
  ↓ Conv1D(1→32) + BatchNorm + ReLU + MaxPool(2)
  ↓ Conv1D(32→64) + BatchNorm + ReLU + MaxPool(2)
  ↓ Flatten
  ↓ Linear(64*5 → 128) + ReLU + Dropout(0.5)
  ↓ Linear(128 → 2) [Softmax]
Output: [P(normal), P(attack)]
\end{verbatim}

\subsubsection{Improved: CNN-LSTM Classifier}

\begin{verbatim}
Input (batch, 20 features)
  ↓ Conv layers (same as baseline)
  ↓ Permute → (batch, seq_len, channels)
  ↓ LSTM(input=64, hidden=64, layers=2, dropout=0.3)
  ↓ Extract last hidden state
  ↓ Linear(64 → 128) + ReLU + Dropout(0.5)
  ↓ Linear(128 → 2) [Softmax]
Output: [P(normal), P(attack)]
\end{verbatim}

\textbf{Rationale:} LSTM layers capture temporal dependencies in network flow sequences, improving detection of multi-step attack patterns.

\subsection{Privacy Mechanisms}

\subsubsection{Baseline: Fixed Differential Privacy}

Following \citet{saklani2026privacy}:

\begin{algorithm}
\caption{Fixed DP-SGD}
\begin{algorithmic}
\STATE \textbf{Input:} Model $\theta$, gradients $g$, clip $C=1.0$, $\epsilon=1.0$, $\delta=10^{-5}$
\STATE \textbf{Compute:} $\sigma = C \sqrt{2\ln(1.25/\delta)} / \epsilon$
\STATE $g \leftarrow g / \max(1, \|g\|_2 / C)$ \COMMENT{Clip}
\STATE $g \leftarrow g + \mathcal{N}(0, \sigma^2 C^2 \mathbf{I})$ \COMMENT{Add noise}
\RETURN $g$
\end{algorithmic}
\end{algorithm}

\subsubsection{Improved: Adaptive Differential Privacy}

Calculates client-specific $\epsilon$ based on data quality:

\begin{algorithm}
\caption{Adaptive Privacy Budget}
\begin{algorithmic}
\STATE \textbf{Input:} Base $\epsilon_0=1.0$, client data quality $q \in [0,1]$
\STATE $\epsilon_{\text{client}} \leftarrow \epsilon_0 \cdot (1 / (q + 0.1))$
\STATE $\epsilon_{\text{client}} \leftarrow \text{clip}(\epsilon_{\text{client}}, 0.5\epsilon_0, 2.0\epsilon_0)$
\RETURN $\epsilon_{\text{client}}$
\end{algorithmic}
\end{algorithm}

Data quality $q$ is estimated via class balance entropy:
\begin{equation}
q = \frac{H(p)}{H_{\text{max}}} = \frac{-\sum_i p_i \log_2 p_i}{\log_2 K}
\end{equation}
where $p_i$ is the proportion of class $i$ and $K$ is the total number of classes.

\subsection{Communication-Efficient Aggregation}

\subsubsection{Top-k Gradient Compression}

\begin{algorithm}
\caption{Top-k Sparsification}
\begin{algorithmic}
\STATE \textbf{Input:} Gradient $g$, compression ratio $r=0.5$
\STATE $k \leftarrow \lfloor r \cdot |g| \rfloor$
\STATE $\text{indices} \leftarrow \text{top-k-by-magnitude}(g, k)$
\STATE $g_{\text{sparse}} \leftarrow \text{zeros-like}(g)$
\STATE $g_{\text{sparse}}[\text{indices}] \leftarrow g[\text{indices}]$
\RETURN $g_{\text{sparse}}$
\end{algorithmic}
\end{algorithm}

Only the top 50\% of gradients by absolute value are transmitted, reducing bandwidth by approximately 45\% while retaining the most informative parameter updates.

\subsubsection{Quality-Weighted FedAvg}

Standard FedAvg weights clients by data size:
\begin{equation}
\theta_{\text{global}} = \sum_{k=1}^{K} \frac{n_k}{n} \theta_k
\end{equation}

SecureFL-IDS incorporates data quality:
\begin{equation}
\theta_{\text{global}} = \sum_{k=1}^{K} \frac{q_k \cdot n_k}{\sum_j q_j \cdot n_j} \theta_k
\end{equation}

where $q_k$ is client $k$'s data quality score, down-weighting contributions from clients with highly imbalanced or low-quality data.

\subsection{Federated Learning Protocol}

\begin{algorithm}
\caption{SecureFL-IDS Training Protocol}
\begin{algorithmic}
\STATE \textbf{Server initializes} global model $\theta_0$
\FOR{round $t = 1$ to $T$}
    \STATE \textbf{Server} broadcasts $\theta_t$ to all clients
    \FOR{each client $k \in \{1, \ldots, K\}$ in parallel}
        \STATE $\epsilon_k \leftarrow \text{AdaptivePrivacy}(q_k)$ \COMMENT{Improved only}
        \STATE $\theta_k \leftarrow$ LocalTrain($\theta_t$, $\mathcal{D}_k$, $\epsilon_k$)
        \STATE $\Delta_k \leftarrow \theta_k - \theta_t$
        \STATE $\Delta_k \leftarrow$ Compress($\Delta_k$) \COMMENT{Improved only}
        \STATE \textbf{Send} $\Delta_k$ to server
    \ENDFOR
    \STATE $\theta_{t+1} \leftarrow \theta_t + \text{Aggregate}(\{\Delta_k\})$
    \STATE \textbf{Evaluate} $\theta_{t+1}$ on test set
\ENDFOR
\RETURN final model $\theta_T$
\end{algorithmic}
\end{algorithm}

\subsection{Deployment Considerations}

\subsubsection{Local Simulation}

Default mode for development and evaluation:
\begin{itemize}
    \item Multi-process clients sharing test set via memory
    \item No network overhead
    \item Deterministic for reproducibility
\end{itemize}

\subsubsection{Containerized Deployment}

Docker Compose configuration:
\begin{itemize}
    \item Separate container per client + server
    \item gRPC or REST API for parameter exchange
    \item Persistent volume for model checkpoints
\end{itemize}

\subsubsection{Cloud Deployment}

AWS ECS or Kubernetes:
\begin{itemize}
    \item Clients as ECS tasks / K8s pods
    \item Server as central service with load balancer
    \item S3 / persistent volumes for model storage
    \item CloudWatch / Prometheus for monitoring
\end{itemize}

\subsection{Security Considerations}

\textbf{Threat Model Assumptions:}
\begin{itemize}
    \item Honest-but-curious server (follows protocol but may infer from updates)
    \item No Byzantine clients (all follow training protocol)
    \item No model poisoning or gradient manipulation attacks
\end{itemize}

\textbf{Privacy Guarantees:}
\begin{itemize}
    \item ($\epsilon$, $\delta$)-DP protects against membership inference
    \item Gradient clipping bounds sensitivity to individual records
    \item No raw data leaves client boundaries
\end{itemize}

\textbf{Limitations:}
\begin{itemize}
    \item Not robust to Byzantine attacks (future work: Krum/Trimmed Mean)
    \item Vulnerable to sophisticated gradient attacks (future: secure aggregation)
    \item Assumes secure communication channels (implement TLS in production)
\end{itemize}
